% !TeX program = pdflatex
% Handout: Unfallatlas Germany (QUA³CK Q→K), Projektpräsentation.
\documentclass[11pt,a4paper]{article}

% ---------- Sprache, Schrift, Mikrotypografie ----------
\usepackage[a4paper,margin=2.1cm,top=2.2cm,bottom=2.2cm]{geometry}
\usepackage[ngerman]{babel}
\usepackage{microtype}
\tolerance=9999
\setlength{\emergencystretch}{3em}
\hyphenpenalty=200
\exhyphenpenalty=50
\usepackage[T1]{fontenc}
\usepackage[utf8]{inputenc}
\usepackage[scale=0.94]{tgheros}
\renewcommand{\familydefault}{\sfdefault}

% ---------- Farbschema (Projekt-Branding, siehe README/streamlit_app.py) ----------
\usepackage{xcolor}
\definecolor{Accent}{HTML}{315F7D}   % steel blue - Primärfarbe der App/Notebooks
\definecolor{AccentDark}{HTML}{24475F}
\definecolor{KSI}{HTML}{E63946}      % Getötet/Schwerverletzt (severity encoding)
\definecolor{Slight}{HTML}{2A9D8F}   % Leichtverletzt
\definecolor{TextGray}{HTML}{2E2E2E}
\definecolor{BoxBg}{HTML}{F1F5F8}
\color{TextGray}

% ---------- Absätze, Listen ----------
\setlength{\parindent}{0pt}
\setlength{\parskip}{4pt}
\usepackage{enumitem}
\setlist[itemize]{nosep,leftmargin=1.2em,topsep=2pt,itemsep=2pt,parsep=0pt}

% ---------- Tabellen ----------
\usepackage{booktabs}
\usepackage{tabularx}
\usepackage{array}
\newcolumntype{Y}{>{\raggedright\arraybackslash}X}
\renewcommand{\arraystretch}{1.15}

% ---------- Boxen ----------
\usepackage[skins]{tcolorbox}
\tcbset{
  basebox/.style={
    colback=BoxBg, colframe=Accent, boxrule=0.5pt, arc=1.5pt,
    left=6pt, right=6pt, top=4pt, bottom=4pt, boxsep=1pt,
    fontupper=\small
  }
}

% ---------- Überschriften ----------
\usepackage{titlesec}
\titleformat{\section}
  {\normalfont\Large\bfseries\color{Accent}}
  {\thesection}{0.5em}{}
  [{\color{Accent}\titlerule[1.1pt]}]
\titlespacing*{\section}{0pt}{12pt}{6pt}
\titleformat{\subsection}
  {\normalfont\large\bfseries\color{AccentDark}}
  {}{0em}{}
\titlespacing*{\subsection}{0pt}{8pt}{3pt}

% ---------- Kopf-/Fußzeile ----------
\usepackage{fancyhdr}
\pagestyle{fancy}
\fancyhf{}
\renewcommand{\headrulewidth}{0.4pt}
\renewcommand{\headrule}{\hbox to\headwidth{\color{Accent}\leaders\hrule height \headrulewidth\hfill}}
\fancyhead[L]{\footnotesize\color{Accent}\bfseries Unfallatlas Germany: QUA\textsuperscript{3}CK}
\fancyhead[R]{\footnotesize\color{Accent} Jonas Weirauch}
\fancyfoot[C]{\footnotesize\thepage\,/\,\pageref{LastPage}}

% ---------- Links, Symbole ----------
\usepackage{amssymb}
\usepackage[hidelinks,colorlinks=true,linkcolor=Accent,urlcolor=Accent]{hyperref}
\usepackage{lastpage}

% ---------- Kleine Helfer ----------
\newcommand{\code}[1]{\texttt{\small #1}}

\begin{document}
\thispagestyle{empty}

\begin{center}
  {\Huge\bfseries\color{Accent} Unfallatlas Deutschland}\par\vspace{4pt}
  {\Large\bfseries Schweregrad-Klassifikation von Verkehrsunfällen in Deutschland}\par\vspace{6pt}
  {\small Projekt-Handout nach dem QUA\textsuperscript{3}CK-Prozessmodell\\
  Question~\textrightarrow{}~Understanding the Data~\textrightarrow{}~Algorithm/Adapt/Adjust~\textrightarrow{}~Conclude \& Compare~\textrightarrow{}~Knowledge Transfer\par\vspace{3pt}
  Jonas Weirauch \textbar{} IU Internationale Hochschule, B.Sc. Data Analytics / Big Data\par
  \href{https://github.com/jonasyr/unfallatlas-qua3ck}{github.com/jonasyr/unfallatlas-qua3ck} \textbar{}
  \href{https://jonasyr.github.io/unfallatlas-qua3ck/}{Live-Report (GitHub Pages)} \textbar{}
  \href{https://unfallatlas-qua3ck-fg5thv4lthkbhw86er3kfk.streamlit.app}{Streamlit-App}}
\end{center}

\begin{tcolorbox}[basebox]
\textbf{Kurzfassung:} Aus 2{,}09~Mio. Unfallatlas-Datensätzen (2016-2024) plus offen
lizenzierten Wetter- (DWD) und Straßennetz-Daten (OSM) wird die
Unfallschwere vorhergesagt. Das ursprüngliche 3-Klassen-Ziel erweist
sich als mit den verfügbaren Merkmalen nicht erreichbar
(Ceiling Macro-F1~$=0{,}424$) und wird auf eine binäre KSI-Variable
(getötet/schwerverletzt vs.\ leichtverletzt) reframt. Der finale
Random-Forest-Champion erreicht auf dem unberührten Testjahr~2024
Macro-F1~$=0{,}604$ und Recall(KSI)~$=0{,}515$, beide Zielkriterien
erfüllt. Das Ergebnis ist als interaktive Streamlit-App und als
statische Notebook-Präsentation live nutzbar.
\end{tcolorbox}

% =====================================================================
\section{Q: Frage}
\textbf{Forschungsfrage:} Welche raumzeitlichen, infrastrukturellen und
meteorologischen Faktoren bestimmen die Schwere eines Verkehrsunfalls in
Deutschland, und lässt sich diese Schwere zuverlässig aus öffentlich
verfügbaren Daten mittels interpretierbarer Machine-Learning-Modelle
vorhersagen?

\textbf{Hypothesen:} (H1) Unfallzeitpunkt (Wochentag, Stunde, Monat)
korreliert mit der Schwere, z.\,B.\ höhere Schwere bei Nachtfahrten
oder außerorts. (H2) Beteiligungsart (Fußgänger, Rad, Motorrad) ist ein
stärkerer Schwere-Prädiktor als reine Ortsmerkmale, da ungeschützte
Verkehrsteilnehmer überproportional betroffen sind. (H3) Wetter- und
Straßeninfrastrukturmerkmale liefern einen messbaren, aber
untergeordneten Zusatzbeitrag gegenüber Unfallart/-typ.

\textbf{Datenquelle:} Unfallatlas Deutschland (GovData/Mobilithek,
Statistisches Bundesamt), Datenlizenz Deutschland, Namensnennung~2.0,
2016-2024, rund 2{,}09~Mio. polizeilich erfasste Unfälle mit
Personenschaden. Zielvariable \code{UKATGEORIE}: 1~=~Getötet~(1\,\%),
2~=~Schwerverletzt~(18\,\%), 3~=~Leichtverletzt~(81\,\%): eine stark
unbalancierte 3-Klassen-Variable.

\textbf{Teststrategie:} Chronologischer Split statt Zufallssplit, da die
Daten Zeitreihencharakter haben (Unfallzahlen, Fahrzeugbestand und
Infrastruktur verändern sich über die Jahre): Train 2016-2022,
Validierung 2023, Test 2024. Damit wird das Modell stets an
\emph{zukünftigen}, nie an zufällig gemischten Daten geprüft, ein
realistischeres Szenario für ein produktiv genutztes Frühwarn- oder
Priorisierungswerkzeug als ein zufälliger Train/Test-Split.

\textbf{Relevanz \& Einordnung:} Interpretierbare Schwere-Vorhersage
ist für Verkehrssicherheitsplanung (Ressourcenpriorisierung bei
Unfallhäufungsstellen), Versicherungswesen und Rettungsdienst-Planung
relevant. Die bewusste Wahl klassischer, gut erklärbarer Modelle
(Random~Forest, Boosting, lineare Modelle) statt Deep-Learning-Ansätzen
folgt der Anforderung, dass Entscheidungsträger:innen die
Vorhersagegründe nachvollziehen können müssen (siehe C-Phase,
Permutation Importance).

\begin{center}
\small
\begin{tabularx}{\linewidth}{Yr}
\toprule
\textbf{Zielmetrik} & \textbf{Schwelle} \\
\midrule
Macro-F1 (Test 2024) & $\geq 0{,}55$ \\
Recall Klasse~1 (Getötet) & $\geq 0{,}50$ \\
Baseline Macro-F1 (Mehrheitsklasse) & $\approx 0{,}30$ \\
\bottomrule
\end{tabularx}
\end{center}

\clearpage

% =====================================================================
\section{U: Untersuchung der Daten}
Explorative Datenanalyse (EDA) und Feature-Engineering ergänzen die
21~Rohspalten des Unfallatlas um externe, offen lizenzierte Quellen,
alles offline reproduzierbar über \code{duckdb} auf dem konsolidierten
Parquet (\code{data/accidents.parquet}, 65~MB, verwaltet via Git~LFS),
ohne die 2{,}09~Mio. Zeilen vollständig nach pandas laden zu müssen.

\begin{itemize}
\item \textbf{Wetter (DWD, CDC Open Data):} Windgeschwindigkeit,
  Niederschlag; je Unfall wird die räumlich/zeitlich nächste
  DWD-Station verknüpft (\code{dwd\_station\_id}, \code{dwd\_station\_dist\_km}).
\item \textbf{Straßennetz (OpenStreetMap, ODbL):} Straßendichte und
  zulässige Höchstgeschwindigkeit im Umkreis, aggregiert über
  H3-Hexagon-Zellen statt teurer Punkt-für-Punkt-Lookups gegen das
  OSM-Straßennetz: notwendig, um den Join auf 2{,}09~Mio. Punkten in
  vertretbarer Zeit auf einem normalen Laptop durchzuführen.
\item \textbf{Zeitliche Merkmale:} Wochentag (\code{UWOCHENTAG}),
  Stunde (\code{USTUNDE}), Monat (\code{UMONAT}): die EDA zeigt klare
  Spitzen im Berufsverkehr sowie saisonale Schwankungen.
\item \textbf{Räumliche Merkmale:} Regierungsbezirk-/Kreis-Code
  (\code{UREGBEZ}/\code{UKREIS}), Breiten-/Längengrad. Ohne
  Bundesland-Schlüssel (\code{ULAND}) im Inference-Contract sind
  (\code{UREGBEZ},\code{UKREIS})-Paare \emph{nicht} global eindeutig auf
  einen amtlichen Kreisnamen abbildbar; die App zeigt sie deshalb
  bewusst als dataset-interne Codes statt erfundener Kreisnamen.
\item \textbf{Beteiligungsmerkmale:} Binärflags je Verkehrsmittel
  (\code{IstPKW}, \code{IstRad}, \code{IstKrad}, \code{IstFuss},
  \code{IstGkfz}, \code{IstSonstig}).
\item \textbf{Unfalltyp \& Zustand:} Unfallart/-typ (\code{UART},
  \code{UTYP1}), Lichtverhältnisse (\code{ULICHTVERH}) und
  Straßenzustand (\code{STRZUSTAND}).
\end{itemize}

\textbf{Zentraler EDA-Befund:} Die Zielverteilung ist extrem
unbalanciert (1\,\% Getötet, 81\,\% Leichtverletzt), und einzelne
Merkmale korrelieren nur schwach mit der Schwere (Cramér's~V
$\leq 0{,}13$ für die stärksten Merkmale). Diese schwache
Merkmals-Ziel-Assoziation ist der entscheidende Vorbote für den
Ceiling-Befund der A\textsuperscript{3}-Phase.

\textbf{Datenqualität:} Vor der Modellierung wurden Plausibilitätsprüfungen
durchgeführt (Wertebereiche der Koordinaten gegen die Landesgrenzen
Deutschlands, Konsistenz der Kategorie-Codes gegen die amtliche
Datensatzbeschreibung \code{DSB\_Unfallatlas.md}, Duplikatsprüfung über
die Unfall-ID). Fehlende DWD-/OSM-Zuordnungen (z.\,B.\ bei sehr
abgelegenen Unfallorten) werden nicht imputiert, sondern als fehlender
Wert weitergereicht, damit Modelle sie explizit als Information nutzen
können, statt eine künstliche Zuordnung vorzutäuschen.

% \clearpage

% =====================================================================
\section{A\textsuperscript{3}: Algorithmus, Anwendung, Anpassung}
\subsection*{Der 3-Klassen-Ceiling und die Reframing-Entscheidung}
Über 19~getestete Konfigurationen, darunter Baselines, lineare Modelle,
Baum-/Boosting-Modelle, Imbalance-Strategien (Klassengewichtung,
Over-/Undersampling) sowie Optuna-basiertes Hyperparameter-Tuning,
erreicht das 3-Klassen-Ziel nur einen empirischen \textbf{Ceiling von
Macro-F1 $= 0{,}424$}, deutlich unter dem Akzeptanzkriterium von
$0{,}55$. Für eine brauchbare Precision der Getötet-Klasse wäre eine
$\sim\!90$-fache Odds-Lift nötig, ein Effekt, den die verfügbaren
Merkmale (Cramér's~V $\leq 0{,}13$) schlicht nicht hergeben.

\textbf{Konsequenz:} Reframing auf eine \textbf{binäre KSI-Variable}
(Killed~or~Seriously~Injured vs.\ Leichtverletzt), die in der
internationalen Unfallforschungs-Literatur etablierte
Standardformulierung (Santos 2022; Pakgohar 2021; Schloessler 2024),
die das seltene, sicherheitskritische Ereignis von der
Alltagsverletzung trennt, statt drei ungleich seltene Klassen simultan
zu modellieren.

\subsection*{Kandidatenvergleich (Stufe 0/1, Validierung 2023)}
Zehn klassengewichtete/balancierte Kandidaten wurden auf identischer
Merkmalsbasis verglichen; die Auswahl erfolgt gate-bewusst
(\code{select\_best\_candidate}): höchster Macro-F1 unter allen
Kandidaten mit Recall(KSI)~$\geq 0{,}50$, sonst höchster Kombi-Score.

\begin{center}
\small
\begin{tabularx}{\linewidth}{Yrrr}
\toprule
\textbf{Modell} & \textbf{Macro-F1} & \textbf{Rec.\,KSI} & \textbf{Rec.\,Leicht} \\
\midrule
Zufallsrat. & 0{,}439 & 0{,}499 & 0{,}499 \\
Mehrheitsklasse & 0{,}453 & 0{,}000 & 1{,}000 \\
Logistic Regression & 0{,}558 & 0{,}638 & 0{,}634 \\
\textbf{Random Forest (Champion)} & \textbf{0{,}601} & \textbf{0{,}540} & \textbf{0{,}747} \\
XGBoost & 0{,}570 & 0{,}682 & 0{,}633 \\
LightGBM & 0{,}566 & 0{,}690 & 0{,}624 \\
CatBoost & 0{,}565 & 0{,}686 & 0{,}624 \\
SVM (linear) & 0{,}559 & 0{,}636 & 0{,}636 \\
SVM (RBF) & 0{,}560 & 0{,}606 & 0{,}651 \\
SVM (SGD) & 0{,}553 & 0{,}654 & 0{,}618 \\
\bottomrule
\end{tabularx}
\end{center}

\textbf{Champion:} Random Forest (balanciert), Hyperparameter
\code{n\_estimators=180}, \code{max\_depth=23},
\code{min\_samples\_leaf=8}; Entscheidungsschwelle $0{,}4986$
(auf Val~2023 optimiert). Val-Macro-F1 $0{,}608$, Recall(KSI)
$0{,}505$: beide Zielgates auf Val~2023 erfüllt. Trainingsumfang:
$\sim$1{,}55~Mio. Zeilen (2016-2022); SVM-Varianten wurden aus
Laufzeitgründen auf Subsamples (8\,000-500\,000~Zeilen) trainiert.

% \clearpage

% =====================================================================
\section{C: Auswertung \& Vergleich}
\subsection*{Testergebnis (Held-out 2024)}
Macro-F1 $= 0{,}604$, Recall(KSI) $= 0{,}515$, Recall(Leicht)
$= 0{,}769$: beide Akzeptanzgates werden auf den unberührten
2024-Daten erneut erfüllt, ein wichtiges Indiz gegen Overfitting auf
den Validierungssatz.

\subsection*{Warum der Champion nicht der Top-Kandidat ist}
Auf der \emph{kombinierten} Kriterienmenge (Recall, Inferenzlatenz,
Robustheit über Resampling-Wiederholungen) liegen die drei
Boosting-Finalisten (XGBoost/LightGBM/CatBoost) beim Recall(KSI)
klar vorne ($0{,}68$-$0{,}69$ vs.\ $0{,}54$) und sind gleichzeitig
$3$-$4\times$ schneller in der Inferenz (12-15~ms/1000~vs.\
43{,}6~ms/1000~Zeilen). Nach diesen Zahlen ist Random~Forest also
\emph{nicht} der bevorzugte Finalist. Er bleibt dennoch
\textbf{Deployment-Champion}, weil die Modellauswahl bereits mittels
Validierungsdaten festgelegt wurde, \emph{bevor} der Testsatz
angefasst wurde: der Held-out-2024-Satz kann nicht nachträglich zur
Neuauswahl zwischen Kandidaten verwendet werden, ohne die
statistische Unabhängigkeit der bereits für Random~Forest berichteten
Testmetriken zu verletzen. Dieser methodische Punkt (Validierung
entscheidet, der Test bestätigt nur) ist zentral für die
Glaubwürdigkeit des berichteten Ergebnisses und wird in der App
(Modellvergleich-Seite) explizit erläutert.

\subsection*{Permutation Importance (Champion-Modell)}
SHAP wurde zugunsten von Permutation Importance verworfen (bessere
Skalierbarkeit auf $\sim$270\,000 Testzeilen ohne Kernel-Approximation).
Top-Merkmale:

\begin{center}
\small
\begin{tabularx}{\linewidth}{Yr}
\toprule
\textbf{Merkmal} & \textbf{Importance} \\
\midrule
\code{UART} (Unfallart) & 0{,}0304 \\
\code{UTYP1} (Unfalltyp) & 0{,}0175 \\
\code{IstKrad} (Motorrad beteiligt) & 0{,}0129 \\
\code{IstRad} (Fahrrad beteiligt) & 0{,}0106 \\
\code{STRZUSTAND} (Straßenzustand) & 0{,}0082 \\
\code{USTUNDE} (Unfallstunde) & 0{,}0066 \\
\bottomrule
\end{tabularx}
\end{center}
Unfallart/-typ und die Beteiligung ungeschützter Verkehrsteilnehmer
(Motorrad, Fahrrad) dominieren die Vorhersage, plausibel und
konsistent mit Hypothese~H2, aber \emph{korrelativ, nicht kausal}.

\subsection*{Limitationen}
Polizeiliche Meldequote variiert nach Schweregrad (schwerere Unfälle
werden zuverlässiger erfasst als leichte, sogenannter Meldebias);
Regierungsbezirk/Kreis-Codes sind ohne Bundesland-Schlüssel nicht
eindeutig lokalisierbar; Rasterzellen der Kartenansicht sind keine
Verwaltungsgrenzen; relative Risikowerte in der Overview-Karte sind
statistische Assoziation, kein Kausalnachweis; Shrinkage-Glättung
($k=20$) verzerrt kleine Stichproben absichtlich Richtung
Landesdurchschnitt, um Overfitting auf dünn besetzte Zellen zu
vermeiden, ein bewusster Fairness-/Robustheits-Trade-off, kein Fehler.

% \clearpage

% =====================================================================
\section{K: Wissenstransfer (Kommunikation)}
Die Ergebnisse werden über zwei sich gegenseitig verlinkende
Oberflächen kommuniziert: eine interaktive Streamlit-App und eine
statische Notebook-Präsentation.

\subsection*{Streamlit-App (\code{app/streamlit\_app.py})}
Läuft vollständig aus committeten \code{data/processed/}-Artefakten
(Modelle, Contract, Kartendaten); kein Notebook-Rerun nötig.

\begin{itemize}
\item \textbf{Overview:} Kennzahlen-Konsole (Macro-F1, Recall(KSI),
  Entscheidungsschwelle), Gegenüberstellung 3-Klassen-Ceiling vs.\
  Binär-Ergebnis, interaktive OSM-Karte mit KSI-Risiko relativ zum
  Landesdurchschnitt (diskrete Risikobänder, einzeln umschaltbar,
  Konfidenz-Transparenz für dünn besetzte Zellen).
\item \textbf{Risikoprädiktor:} Klickbare Karte (Folium/OSM) statt
  manueller Lon/Lat-Eingabe: ein Klick setzt Position und Formular
  in einem Schritt; Formular für Kontextmerkmale (Zeit, Beteiligung,
  Straßenzustand); Live-Vorhersage inkl.\ Konfidenz relativ zur
  Entscheidungsschwelle.
\item \textbf{Warum diese Vorhersage:} Permutation-Importance-Chart
  und Eingabetabelle mit menschenlesbaren Feature-Labels (z.\,B.
  \code{dwd\_wind\_speed\_ms}~\textrightarrow{}~\glqq{}Windgeschwindigkeit
  (m/s)\grqq{}) statt Rohspaltennamen.
\item \textbf{Modellvergleich:} Alle 10~Kandidaten, Pareto-Front
  (Macro-F1 vs.\ Recall(KSI)), Konfusionsmatrix des Champions
  (Test 2024), Erklärung der Champion-vs-Finalist-Entscheidung
  (siehe C-Phase oben).
\end{itemize}

\subsection*{Notebook-Präsentation \& Deployment}
Die vier ausgeführten Q/U/A\textsuperscript{3}/C-Notebooks werden
zusätzlich als statisches, offline-fähiges HTML (Code, Text, sowie
interaktive Plotly-Grafiken) über GitHub~Pages veröffentlicht, ohne
lokale Installation einsehbar. App und Notebook-Präsentation sind
gegenseitig verlinkt, sodass Leser:innen zwischen
Analyse-Nachvollziehbarkeit (Notebooks) und interaktiver Exploration
(App) wechseln können. Deployment der App erfolgt über Streamlit
Community Cloud, Branch~\code{main}, ohne dass Notebooks zur Laufzeit
neu ausgeführt werden müssen.

\subsection*{Fazit \& Ausblick}
Das Projekt zeigt exemplarisch, wie ein zunächst nicht erreichbares
Klassifikationsziel durch begründetes, literaturgestütztes Reframing
(3-Klassen~\textrightarrow{}~binäre KSI-Variable) doch noch belastbare,
handlungsrelevante Ergebnisse liefert, und wie wichtig die saubere
Trennung von Validierungs- und Testdaten für die Glaubwürdigkeit einer
Modellauswahl ist. Mögliche Erweiterungen: Einbindung des amtlichen
Gemeindeschlüssels (\code{ULAND}) für belastbare Kreis-Namen,
Kalibrierung der Vorhersagewahrscheinlichkeiten, sowie eine
zeitbasierte \glqq{}Risk-Scan\grqq{}-Simulation auf der Overview-Karte
(aktuell bewusst zurückgestellt, um Umfang und Qualität der
bestehenden Fixes nicht zu gefährden).

\subsection*{Tech-Stack}
\begin{center}
\small
\begin{tabularx}{\linewidth}{lY}
\toprule
\textbf{Bereich} & \textbf{Bibliotheken} \\
\midrule
Daten & pandas, polars, duckdb, pyarrow \\
ML & scikit-learn, xgboost, lightgbm, catboost \\
Imbalance \& Tuning & imbalanced-learn, optuna \\
Erklärbarkeit & Permutation Importance (statt SHAP) \\
Geodaten & folium, streamlit-folium (optional: geopandas, h3, osmnx) \\
App & streamlit \\
Tooling & uv, ruff, pytest, jupytext, pre-commit \\
\bottomrule
\end{tabularx}
\end{center}

\subsection*{Projektstruktur \& Reproduzierbarkeit}
Der gesamte Ablauf ist ohne manuelle Zwischenschritte reproduzierbar:
\code{git lfs pull} (Datensatz) $\rightarrow$ \code{uv sync -{}-all-extras}
(Abhängigkeiten) $\rightarrow$ \code{uv run streamlit run
app/streamlit\_app.py} (App) bzw.\ \code{uv run jupyter lab}
(Notebooks). Kein Schritt erfordert einen Notebook-Rerun, da alle für
App und Präsentation nötigen Artefakte bereits unter
\code{data/processed/} committet sind.

\begin{center}
\small
\begin{tabularx}{\linewidth}{lY}
\toprule
\textbf{Verzeichnis} & \textbf{Inhalt} \\
\midrule
\code{notebooks/} & QUA\textsuperscript{3}CK-Phasennotebooks (Source of Truth): Q, U, A\textsuperscript{3}, C \\
\code{src/unfallatlas/} & Wiederverwendbare Bibliothek (\code{data/}, \code{features/}, \code{models/}, \code{viz/}, \code{presentation/}) \\
\code{app/} & Streamlit-Demo (K-Phase): Overview, Risikoprädiktor, Warum diese Vorhersage, Modellvergleich \\
\code{data/} & \code{accidents.parquet} (Git~LFS) sowie \code{interim/}- und \code{processed/}-Artefakte (Modelle, Contract, Kartendaten) \\
\code{tests/} & pytest-Suite inkl.\ Coverage-Report und optionalen Playwright-Browsertests \\
\code{docs/} & Disclosure, Glossar, Datensatzbeschreibung, Prompt-Transkripte je Phase \\
\code{reports/} & Generierte Abbildungen und Notebook-Präsentationen (GitHub~Pages) \\
\bottomrule
\end{tabularx}
\end{center}

Qualitätssicherung läuft identisch lokal und in CI (GitHub Actions):
\code{uv run pytest} (Tests inkl.\ Coverage), \code{uv run ruff check .}
(Linting), \code{uv run ruff format .} (Formatierung) sowie
Pre-Commit-Hooks (Ruff, nbstripout, Commitizen, Notebook-Mirror-Check).

\vspace{2pt}
\noindent\rule{\linewidth}{0.4pt}\par\vspace{2pt}
{\small\textbf{Weiterführende Ressourcen:}
\href{https://github.com/jonasyr/unfallatlas-qua3ck}{GitHub-Repository} \textbar{}
\href{https://jonasyr.github.io/unfallatlas-qua3ck/}{Live-Report (Notebooks als HTML)} \textbar{}
\href{https://unfallatlas-qua3ck-fg5thv4lthkbhw86er3kfk.streamlit.app}{Streamlit-App} \textbar{}
Dokumentation: \code{docs/GLOSSARY.md}, \code{docs/AI TOOL DISCLOSURE.md}, \code{docs/dataset/}\par}

{\scriptsize Code: MIT-Lizenz. Unfalldaten: Datenlizenz Deutschland,
Namensnennung 2.0 (Mobilithek/Statistisches Bundesamt, Unfallatlas
Deutschland). Wetterdaten: Deutscher Wetterdienst (DWD), CDC Open Data.
Straßennetz: \copyright{} OpenStreetMap-Mitwirkende, ODbL.}

\end{document}
